\chapter{Experimental Design}
\label{chap:experimental_design}

This chapter should explain \emph{how} the thesis evaluates the proposed rolling-horizon delivery framework, \emph{why} the retained experiment scope is sufficient for the thesis claims, and \emph{how} reproducibility and auditability are maintained. The chapter should read as the bridge between the methodology chapters and the results chapter.

\section{Chapter Purpose and Writing Strategy}

The purpose of this chapter is fourfold:
\begin{enumerate}
    \item define the retained experiment scope of the thesis;
    \item describe the datasets, scenarios, and evaluation pipeline;
    \item justify the choice of baselines, controller progression, and performance metrics;
    \item document the reproducibility and audit process.
\end{enumerate}

The writing tone of this chapter should be precise and restrained. It should avoid overstating causal conclusions and should instead focus on:
\begin{itemize}
    \item what is evaluated,
    \item under which assumptions,
    \item with which controls,
    \item and by which output artifacts.
\end{itemize}

\noindent\textbf{Writing rule for this chapter:}
the main experimental backbone should remain limited to:
\begin{itemize}
    \item \texttt{EXP00},
    \item \texttt{EXP-BASELINE} where needed as implementation context,
    \item \texttt{EXP01},
    \item and the \texttt{EXP01 / Scenario1} controller line:
    \texttt{v2 $\rightarrow$ v4 $\rightarrow$ v5 $\rightarrow$ v6f $\rightarrow$ v6g}.
\end{itemize}

\section{Retained Experiment Scope}

This section should clearly state that the thesis does not attempt to report every historical experiment branch. Instead, it intentionally narrows the evidence base to a retained, reproducible set of experiments.

\subsection{Why the Scope is Narrowed}

Write this subsection as a methodological decision, not as a limitation to hide:
\begin{itemize}
    \item the broader research program explored a richer design space;
    \item however, a graduation thesis requires a stable and reproducible central line;
    \item therefore, the final experimental scope is reduced to the smallest set of experiments that still supports the thesis claims.
\end{itemize}

\noindent Suggested framing:
\begin{quote}
The broader research process considered richer optimization and controller variants, but the final thesis adopts a deliberately retained experiment scope in order to keep the comparison protocol stable, interpretable, and reproducible.
\end{quote}

\subsection{Official Experiment Families}

This subsection should define the meaning of the retained experiments:
\begin{description}
    \item[\texttt{EXP00}] normal-operation reference experiment;
    \item[\texttt{EXP-BASELINE}] implementation baseline, mainly useful for contextualizing the pipeline;
    \item[\texttt{EXP01}] crunch-condition baseline experiment;
    \item[\texttt{Scenario1}] controller comparison line built on top of \texttt{EXP01}.
\end{description}

\noindent The key message is:
\begin{itemize}
    \item \texttt{EXP00} and \texttt{EXP01} define the operational pressure contrast;
    \item \texttt{Scenario1} defines the method-improvement line inside the stressed regime.
\end{itemize}

\section{Data and Instance Construction}

This section should explain the data inputs used by the retained experiments and how they are transformed into runnable simulation instances.

\subsection{Raw Data Source}

State that the source data originate from the operational workbook and that the raw data are treated as read-only. Mention that the raw workbook provides:
\begin{itemize}
    \item orders,
    \item warehouse/depot settings,
    \item vehicle settings.
\end{itemize}

This subsection should briefly summarize the operational semantics already established earlier in the thesis:
\begin{itemize}
    \item release dates and due dates,
    \item feasible service windows,
    \item depot-specific operational timing,
    \item vehicle availability and capacity rules.
\end{itemize}

\subsection{Processed Benchmark and Matrix Inputs}

This subsection should define the retained processed inputs:
\begin{itemize}
    \item processed benchmark instance, with Herlev as the main thesis benchmark;
    \item OSRM-based travel-time and distance matrices;
    \item processed benchmark variants for additional depots where relevant for out-of-distribution validation.
\end{itemize}

\noindent Important wording:
\begin{itemize}
    \item emphasize that the retained thesis line uses road-network matrices rather than Euclidean approximations;
    \item describe processed data as reproducible artifacts derived from the raw source.
\end{itemize}

\subsection{Instance Semantics for the Retained Thesis Line}

This subsection should explain what one experiment instance means in practice:
\begin{itemize}
    \item a rolling-horizon multi-day planning problem,
    \item evaluated day by day,
    \item with only released-but-undelivered orders visible,
    \item and with routing solved subject to same-day operational limits.
\end{itemize}

\section{Scenario Design}

This section should explain the operational conditions under which the experiments are run.

\subsection{Normal and Crunch Regimes}

Define the two main evaluation regimes:
\begin{itemize}
    \item \texttt{EXP00}: business-as-usual reference;
    \item \texttt{EXP01}: stressed single-wave or crunch reference.
\end{itemize}

This subsection should answer:
\begin{itemize}
    \item why the thesis needs both a normal and a stressed regime;
    \item why performance under pressure is the key setting for evaluating controller improvements.
\end{itemize}

\subsection{Scenario1 as the Main Method Comparison Regime}

This subsection should state that \texttt{Scenario1} is not a separate unrelated experiment family. It is the main controller-comparison setting within the retained \texttt{EXP01} stress regime.

The central message should be:
\begin{itemize}
    \item \texttt{EXP01} defines the stressed environment;
    \item \texttt{Scenario1} compares alternative controller designs under that environment;
    \item this makes the controller line directly interpretable as a method progression under pressure.
\end{itemize}

\subsection{Out-of-Distribution Depot Evaluation}

This subsection should explain why selected OOD depot runs are included:
\begin{itemize}
    \item to test whether the final retained controller line is overly tuned to one depot;
    \item to provide a compact generalization check without turning the thesis into a full transfer-learning study.
\end{itemize}

\noindent Important caution:
OOD evidence should be presented as supportive generalization evidence, not as proof of universal robustness.

\section{Baseline and Method Comparison Design}

This section should define which methods are compared and why.

\subsection{Operational Baselines}

Describe the baseline hierarchy:
\begin{enumerate}
    \item normal-operation baseline via \texttt{EXP00},
    \item crunch baseline via \texttt{EXP01},
    \item controller progression within \texttt{Scenario1}.
\end{enumerate}

The baseline logic should be written so that the reader understands:
\begin{itemize}
    \item first, what degradation occurs under stress;
    \item second, how controller upgrades recover performance under that stress.
\end{itemize}

\subsection{Retained Controller Progression}

This subsection should describe the final controller line in prose:
\begin{description}
    \item[\texttt{v2}] early robust baseline;
    \item[\texttt{v4}] event-driven commitment mechanism;
    \item[\texttt{v5}] risk-budgeted admission / commitment / release logic;
    \item[\texttt{v6f}] execution-aware transition step;
    \item[\texttt{v6g}] final deadline-reservation endpoint and main thesis candidate.
\end{description}

This subsection should not yet present the results. It should only explain the intended comparison logic and what each step is meant to test.

\subsection{Excluded or Secondary Variants}

To keep the chapter clean, state explicitly that some variants are not part of the retained main line:
\begin{itemize}
    \item historical branches that do not support the final thesis claim;
    \item counterexample variants such as diversification-focused alternatives;
    \item exploratory follow-up variants that are not required to establish the main thesis result.
\end{itemize}

The goal is to show disciplined scope control.

\section{Execution Protocol}

This section should explain how each experiment is actually run.

\subsection{Rolling-Horizon Evaluation}

Explain the daily loop:
\begin{enumerate}
    \item build the visible set for the current day;
    \item choose the controller action or policy configuration;
    \item solve the day-level routing problem under the computation budget;
    \item commit current-day execution;
    \item roll to the next day and repeat.
\end{enumerate}

This subsection should make clear that evaluation is policy-based, not one-shot static optimization.

\subsection{Computation Budgets}

This subsection should explain the role of bounded compute in the experiment design:
\begin{itemize}
    \item the thesis evaluates methods under finite decision-time budgets;
    \item compute budgets are treated as part of the operational setting;
    \item comparisons must therefore keep compute settings explicit whenever materially relevant.
\end{itemize}

\noindent Include the thesis-safe idea that:
\begin{itemize}
    \item improved performance should not be attributed to physical expansion if it can already be explained by controller logic or compute allocation;
    \item probe experiments should be described carefully and not oversold.
\end{itemize}

\subsection{Fleet and Physical Controls}

This subsection should list the physical or semi-physical controls that must remain explicit during comparison:
\begin{itemize}
    \item maximum trips per vehicle,
    \item active vehicle count,
    \item route duration limits,
    \item compute budget,
    \item fixed scenario definition.
\end{itemize}

The chapter should explain that these controls are held constant unless an experiment explicitly probes them.

\section{Metrics and Evaluation Criteria}

This section should define what the thesis measures and why.

\subsection{Primary Metrics}

List the main outcome metrics:
\begin{itemize}
    \item service rate,
    \item number of failed orders,
    \item penalized cost,
    \item carryover or backlog-related statistics where relevant.
\end{itemize}

Explain that these metrics capture both delivery quality and operational failure.

\subsection{Secondary Metrics}

List secondary or diagnostic metrics:
\begin{itemize}
    \item travel cost and duration,
    \item runtime,
    \item overdue rate,
    \item route-level or execution-level diagnostics,
    \item stability-related measures where available.
\end{itemize}

\subsection{Stability and Interpretation Metrics}

Because the broader research line includes plan stability and operational realism, this subsection should explain:
\begin{itemize}
    \item why stability matters even if it is not always the headline metric;
    \item how stability-aware interpretation remains relevant to the thesis;
    \item how execution-aware analysis helps interpret residual failure gaps.
\end{itemize}

\section{Reproducibility and Audit Trail}

This section should explain how official thesis experiments are kept reproducible.

\subsection{Registered Experiment Definitions}

Describe the role of formal experiment definitions and script-controlled registration. Explain that:
\begin{itemize}
    \item official experiment families are explicitly registered;
    \item ad hoc runs do not define the thesis narrative;
    \item the experiment definition layer provides a stable interface between method description and execution.
\end{itemize}

\subsection{HPC Execution Policy}

State that formal runs are executed through the cluster workflow rather than through casual local execution. This subsection should include:
\begin{itemize}
    \item why official runs use HPC,
    \item why local execution is limited to smoke tests, dry runs, and lightweight validation,
    \item why this distinction matters for experimental discipline.
\end{itemize}

\subsection{Output Artifacts and Audit Files}

Describe the expected outputs of an official experiment:
\begin{itemize}
    \item per-seed summaries,
    \item aggregate statistics,
    \item controller-specific analysis outputs,
    \item logs and audit traces,
    \item figure-building inputs where applicable.
\end{itemize}

The key idea is that the thesis should be traceable from:
\begin{center}
configuration $\rightarrow$ execution $\rightarrow$ outputs $\rightarrow$ tables/figures.
\end{center}

\section{Statistical and Comparative Logic}

This section should explain how comparisons are interpreted.

\subsection{Seeded Comparisons}

State that comparisons are based on controlled repeated runs across seeds where applicable. Emphasize:
\begin{itemize}
    \item seed-matched comparisons are preferred;
    \item reported averages should reflect reproducible aggregates rather than anecdotal best runs.
\end{itemize}

\subsection{Paired Interpretation}

Where the experiment pipeline supports it, explain the use of paired comparisons:
\begin{itemize}
    \item same scenario,
    \item same seed structure,
    \item controlled method change.
\end{itemize}

This subsection should keep the interpretation modest:
\begin{itemize}
    \item the goal is comparative clarity, not overclaimed causal proof.
\end{itemize}

\subsection{Guardrails Against Overclaiming}

This subsection should explicitly connect the experimental design to claim discipline:
\begin{itemize}
    \item baseline differences define pressure effects;
    \item controller differences define method comparisons;
    \item small numerical differences should not automatically be interpreted as strong mechanism proof;
    \item model-side execution limits should not be conflated with real-world physical limits.
\end{itemize}

\section{Threats to Validity}

This section should make the chapter look mature and self-aware.

\subsection{Internal Validity}

Potential issues:
\begin{itemize}
    \item limited retained experiment scope,
    \item dependence on selected stress scenarios,
    \item bounded runtime affecting final routing quality.
\end{itemize}

\subsection{External Validity}

Potential issues:
\begin{itemize}
    \item one main benchmark depot dominates the retained thesis line,
    \item OOD checks are supportive but still limited,
    \item some real operational resources remain abstracted in the model.
\end{itemize}

\subsection{Measurement Validity}

Potential issues:
\begin{itemize}
    \item service-rate improvements may arise from multiple interacting mechanisms;
    \item some secondary metrics are diagnostic rather than direct business outcomes;
    \item absence of a metric in the headline tables does not mean it is irrelevant to interpretation.
\end{itemize}

\section{Chapter-to-Results Handoff}

This short concluding section should explain how the next chapter will present the evidence.

Suggested transition logic:
\begin{enumerate}
    \item first, compare \texttt{EXP00} and \texttt{EXP01} to quantify stress degradation;
    \item second, analyze the retained \texttt{Scenario1} controller progression;
    \item third, identify the retained main candidate \texttt{v6g\_v6d\_compute300};
    \item fourth, use OOD runs as a compact generalization check;
    \item fifth, place \texttt{automode} and the retained \texttt{DATA003} east/west lines in a controlled auxiliary role;
    \item finally, discuss what the evidence does and does not establish.
\end{enumerate}

\section{Writing Checklist for Final Revision}

Before finalizing this chapter, verify that the text does all of the following:
\begin{itemize}
    \item clearly defines \texttt{EXP00}, \texttt{EXP01}, and \texttt{Scenario1};
    \item explains why the scope is narrowed;
    \item makes the controller progression explicit;
    \item names the primary benchmark and matrix assumptions;
    \item defines the primary and secondary metrics;
    \item explains the HPC and reproducibility workflow;
    \item avoids writing results too early;
    \item stays consistent with the retained writing map and claim guardrails.
\end{itemize}

\noindent\textbf{Author note:}
when this chapter is rewritten from plan into final prose, it should remain a methodological chapter. It should set up the evaluation cleanly, but it should save all substantive performance claims, rankings, and final interpretations for the results and discussion chapters.
